% !TeX program = pdflatex
% =============================================================================
% Aufgabenblatt 3: Normalkraft
% UZH Praktikum I -- Sara Romer Urban, Kantonsschule im Lee, HS2026
% Klasse: 2e (09.09.2026)
%
% Reine Uebungsaufgaben zur Normalkraft/schiefen Ebene, kein einleitender
% Erklaerkasten (siehe institutionen/UZH-Praktikum-I/CLAUDE.md). Drei Aufgaben,
% keine Wiederverwendung der Arbeitsblatt-Bilder: (1) Schlitten auf einer neuen
% schiefen Ebene (alpha=20 Grad) -- Gewichtskraft berechnen, dann grafisch in
% Hangabtriebskraft/Normalkraft-Komponente/Normalkraft zerlegen (eigene
% TikZ-Skizze, selbst gewaehlter Kraftmassstab wie beim Lampe-Beispiel aus
% kraefte-addieren-zerlegen), (2) direkte Fortsetzung: aus der grafisch
% bestimmten Hangabtriebskraft Beschleunigung und Endgeschwindigkeit berechnen
% (LZ5, neue Zahlenwerte gegenueber der Bergziege im Arbeitsblatt), (3) kurze
% Ranking-Aufgabe (LZ2, ohne Bild) -- Normalkraft dreier verschieden schwerer
% Fahrzeuge auf waagrechter Strasse vergleichen; bewusst eine andere Perspektive
% auf LZ2 als die Klotz-Aufgabe im Arbeitsblatt (verschiedene Koerper statt
% derselbe Koerper in verschiedenen Situationen), um Redundanz zu vermeiden.
%
% Quellen: lernziele.md, normalkraft-lektionsplan.tex (dieser Ordner).
% =============================================================================
\documentclass[addpoints,11pt,a4paper]{exam}

\usepackage[T1]{fontenc}
\usepackage[utf8]{inputenc}
\usepackage[ngerman]{babel}
\usepackage{lmodern}
\usepackage[a4paper,margin=2.2cm,headheight=14pt]{geometry}
\usepackage{amsmath}
\usepackage{siunitx}
% Hausstil: zusammengesetzte Einheiten immer als echter Bruch (\frac), nicht
% als Inline-Schraegstrich oder \tfrac -- gilt global fuer jedes \si/\SI.
\sisetup{per-mode=fraction,fraction-command=\frac}
\usepackage{array,booktabs,tabularx,multirow}
\usepackage{enumitem}
\usepackage{xcolor}
\usepackage{colortbl}
\usepackage{tikz}
\usepackage{physics}
\usepackage[most]{tcolorbox}
\usepackage{lastpage}
\usepackage{graphicx}
\usepackage{hyperref}

\usetikzlibrary{arrows.meta,calc,angles,quotes,decorations.pathmorphing}

\usepackage{setspace}
\usepackage{needspace}
\setlength{\parindent}{0pt}
\setlength{\parskip}{0.6em}
\setstretch{1.08}
\setlist[itemize]{itemsep=0.4em}

% -----------------------------
% Farben (Corporate-Farbschema Physik KFR)
% -----------------------------
\definecolor{titleblue}{HTML}{183A6B}
\definecolor{accentblue}{HTML}{2E6F9E}
\definecolor{lightblue}{HTML}{EAF4FB}
\definecolor{lightgray}{HTML}{F4F6F8}
\definecolor{darkgray}{HTML}{4A4A4A}
\definecolor{accentorange}{HTML}{D9720C}
\definecolor{accentpurple}{HTML}{7B2D8E}
% Hinweisbox: Orange (accentorange), fuer wichtige Klarstellungen.
\definecolor{lightorange}{HTML}{FCEEE0}
% Formelbox: Violett (accentpurple), fuer die zentralen Merksaetze.
\definecolor{lightpurple}{HTML}{F3EAF7}

% Links (hyperref): farblich ans Corporate-Farbschema angeglichen.
\hypersetup{
  colorlinks=true,
  urlcolor=accentblue,
  linkcolor=accentblue,
  pdftitle={Aufgabenblatt 3: Normalkraft},
  pdfauthor={Loris Keller}
}

% -----------------------------
% Spaltentypen
% -----------------------------
\newcolumntype{Y}{>{\centering\arraybackslash}X}
\newcolumntype{P}[1]{>{\raggedright\arraybackslash}p{#1}}

% -----------------------------
% Boxen
% -----------------------------
\tcbset{before skip=10pt, after skip=10pt}
\newtcolorbox{titlebox}{
  enhanced,
  colback=titleblue,
  colframe=titleblue,
  boxrule=0pt,
  arc=3mm,
  left=3mm,right=3mm,top=3mm,bottom=3mm
}

\newlength{\aufgabenindent}
\settowidth{\aufgabenindent}{10.\hskip\labelsep}

% Praktikum I: Lernziele/Einleitung/Theorie/Aufgaben werden NICHT mehr
% farbig gerahmt (nur Formel-, Hinweis- und Antwortbox bleiben Boxen) --
% siehe institutionen/UZH-Praktikum-I/CLAUDE.md, Abschnitt "Boxen-Layout".
\newenvironment{infobox}{%
  \par\addvspace{10pt}\noindent
}{%
  \par\addvspace{10pt}
}

\newenvironment{theoriebox}[1][Theorie]{%
  \par\addvspace{10pt}\noindent\textbf{#1}\par\smallskip\noindent
}{%
  \par\addvspace{10pt}
}

% taskbox/groupbox stehen in questions VOR dem ersten \question (Bilder,
% Fliesstext) -- exam.cls' questions-Liste braucht dafuer eine echte,
% eigenstaendige Box (sonst "Something's wrong--perhaps a missing \item",
% weil z.B. ein \begin{center} vor dem ersten \item der Aussenliste steht).
% Deshalb bleiben es tcolorboxen, nur unsichtbar gemacht (keine Farbe/Rahmen,
% kein Padding) statt sie durch reine Absaetze zu ersetzen.
\newtcolorbox{taskbox}[1][]{
  enhanced, breakable,
  colback=white, colframe=white, boxrule=0pt,
  left=0mm,right=0mm,top=0mm,bottom=0mm,
  before skip=10pt, after skip=10pt,
  before upper={%
    \def\tmpboxtitle{#1}%
    \ifx\tmpboxtitle\empty\else\noindent\textbf{#1}\par\smallskip\fi
  }
}

\newtcolorbox{groupbox}[1][Gruppenarbeit]{
  enhanced, breakable,
  colback=white, colframe=white, boxrule=0pt,
  left=0mm,right=0mm,top=0mm,bottom=0mm,
  before skip=10pt, after skip=10pt,
  before upper={\noindent\textbf{#1}\par\smallskip}
}

% beispielbox: fuer die Einleitung -- wie die anderen Inhaltsboxen ungerahmt.
\newenvironment{beispielbox}[1][Beispiel]{%
  \par\addvspace{10pt}\noindent\textbf{#1}\par\smallskip\noindent
}{%
  \par\addvspace{10pt}
}

% hinweisbox: Orange, fuer wichtige Klarstellungen/Verwechslungswarnungen.
\newtcolorbox{hinweisbox}[1][Hinweis]{
  enhanced,
  breakable,
  colback=lightorange,
  colframe=accentorange,
  title={#1},
  fonttitle=\bfseries,
  boxrule=0.7pt,
  arc=2mm,
  left=5mm,right=5mm,top=2mm,bottom=2mm
}

% formelbox: Violett, um die zentralen Merksaetze dieser Einheit hervorzuheben.
\newtcolorbox{formelbox}[1][Merksatz]{
  enhanced,
  breakable,
  colback=lightpurple,
  colframe=accentpurple,
  title={#1},
  fonttitle=\bfseries,
  boxrule=0.9pt,
  arc=2mm,
  left=5mm,right=5mm,top=2mm,bottom=2mm
}

\newtcolorbox{answerbox}{
  breakable,
  colback=white,
  colframe=gray!55,
  boxrule=0.5pt,
  arc=1.5mm,
  left=2mm,right=2mm,top=2mm,bottom=2mm
}

% Marke: kleines farbiges Inline-Label fuer Teilschritte/Ergebnis-Callouts.
\newtcbox{\Marke}[1][accentpurple]{
  nobeforeafter,
  colback=#1,
  colframe=#1,
  coltext=white,
  boxrule=0pt,
  arc=2pt,
  left=5pt,right=5pt,top=2pt,bottom=2pt,
  fontupper=\bfseries\sffamily\small
}

% Bild-Helfer: zentriertes Bild mit spuerbarem Abstand zum umgebenden Text
% (statt nur des normalen Absatzabstands). #1 = Breite, #2 = Dateipfad.
\newcommand{\Bild}[2]{%
  \par\vspace{0.6cm}
  \begin{center}
    \includegraphics[width=#1]{#2}
  \end{center}
  \vspace{0.6cm}
}

% --- Loesungen ein-/ausblenden -----------------------------------------------
%\noprintanswers
\printanswers

\pointformat{}
\renewcommand{\solutiontitle}{\noindent\color{titleblue}\textbf{L\"osung:}\enspace}

\pagestyle{headandfoot}
\cfoot{\textcolor{darkgray}{Seite \thepage\ von \pageref{LastPage}}}

\newcommand{\Kopfzeile}[4]{%
  \begin{titlebox}
    {\color{white}\sffamily\bfseries\Large #4}
  \end{titlebox}
  \vspace{0.5em}
  \noindent\textbf{Klasse:}~#1 \hfill \textbf{Datum:}~#2 \hfill \textbf{Thema:}~#3
    \hfill \textbf{Lehrperson:}~Loris Keller
  \vspace{0.8em}
  \lhead{\textcolor{darkgray}{#3}}
  \rhead{\textcolor{darkgray}{#4}}
}

\newlength{\antwortraumlen}
\newcommand{\Antwortraum}[1]{%
  \pgfmathsetlength{\antwortraumlen}{#1*2.5cm}%
  \begin{answerbox}
    \rule{0pt}{\antwortraumlen}%
  \end{answerbox}%
}

% Werte fuer Kopfzeile zentral setzen
\newcommand{\Klasse}{2e}
\newcommand{\Datum}{09.09.2026}
\newcommand{\Thema}{Normalkraft}
\newcommand{\ArbeitsblattTitel}{Aufgabenblatt 3: Normalkraft}

\begin{document}

\Kopfzeile{\Klasse}{\Datum}{\Thema}{\ArbeitsblattTitel}

% =============================================================================
% Aufgabe 1: Schlitten auf schiefer Ebene (Gewichtskraft + grafische Zerlegung)
% =============================================================================
\begin{questions}
\begin{taskbox}[Aufgabe 1: Ein Schlitten auf der Piste]
\question[7] Ein Schlitten mit Kind (Gesamtmasse $m=\SI{30}{\kilogram}$)
steht still auf einer reibungsfreien, verschneiten Rampe mit
Neigungswinkel $\alpha=\SI{20}{\degree}$:

\begin{center}
\begin{tikzpicture}[>=Latex,scale=2.5]
  \def\hangwidth{5}
  \def\winkelgrad{20}
  \pgfmathsetmacro{\hoehe}{\hangwidth*tan(\winkelgrad)}
  \coordinate (F) at (0,0);
  \coordinate (S) at (\hangwidth,\hoehe);
  \draw[thick,darkgray] (F) -- (S);
  \draw[thick,darkgray] (F) -- ++(\hangwidth,0);
  \draw[dashed,gray] (F) -- ++(\hangwidth,0);
  \coordinate (H) at ($(F)+(1,0)$);
  \pic[draw=accentblue,angle radius=0.8cm] {angle=S--F--H};
  \node[accentblue] at ($(F)+(1.05,0.28)$) {$\alpha$};
  \coordinate (K) at ($(F)!0.55!(S)$);
  \draw[fill=lightgray!60,rotate around={\winkelgrad:(K)}]
    ($(K)+(-0.35,0)$) rectangle ++(0.7,0.5);
  \node[rotate=\winkelgrad,inner sep=1pt] at ($(K)+(-{0.42*sin(\winkelgrad)},{0.42*cos(\winkelgrad)})$) {\scriptsize Schlitten};
  \coordinate (SP) at ($(K)+(-{0.25*sin(\winkelgrad)},{0.25*cos(\winkelgrad)})$);
  \fill (SP) circle (1pt);
  \draw[->,ultra thick,accentorange] (SP) -- ++(0,-1.6) node[midway,right]{$\vec F_g$};
\end{tikzpicture}
\end{center}

\begin{parts}
\part[2] Berechnen Sie die Gewichtskraft $F_g$ von Schlitten und Kind
zusammen.

\Antwortraum{1.5}

\begin{solution}
\[
F_g = m\cdot g = \SI{30}{\kilogram}\cdot\SI{9.81}{\meter\per\second\squared}
    \approx\SI{294}{\newton}
\]
\end{solution}

\part[5] Zeichnen Sie $\vec F_g$ mit einem selbst gewählten
Kraftmassstab senkrecht nach unten in die Skizze oben ein. Konstruieren
Sie anschliessend grafisch (Hilfslinien parallel/senkrecht zur Rampe,
Kräfteparallelogramm als Rechteck) die parallele Komponente
$\vec F_{G\parallel}$, die senkrechte Komponente $\vec F_{G\perp}$ und
die Normalkraft $\vec F_N$.

\begin{solution}
Zur Kontrolle (nicht vorrechnen, nur für den Konstruktionsvergleich):
Bei einem flachen Winkel von $\SI{20}{\degree}$ ist $F_{G\parallel}$
deutlich kleiner als $F_g$, $F_N$ dagegen fast so gross wie $F_g$, mit
erwarteter Grössenordnung $F_{G\parallel}\approx\SI{100}{\newton}$ und
$F_N\approx\SI{276}{\newton}$.
\end{solution}
\end{parts}
\end{taskbox}
\end{questions}

% =============================================================================
% Aufgabe 2: Endgeschwindigkeit des Schlittens (LZ5)
% =============================================================================
\newpage
\begin{questions}
\setcounter{question}{1}
\begin{taskbox}[Aufgabe 2: Der Schlitten rutscht los]
\question[8] Der Schlitten aus Aufgabe 1 rutscht aus der Ruhe eine
Hanglänge von $s=\SI{12}{\meter}$ hinunter.

\begin{parts}
\part[3] Welche der beiden Kräfte aus Ihrer Konstruktion in Aufgabe 1
ist für die Beschleunigung entlang des Hangs verantwortlich? Nutzen Sie
diese Kraft und das 2. Newtonsche Gesetz, um die Beschleunigung $a$ zu
berechnen.

\Antwortraum{2}

\begin{solution}
Verantwortlich ist die parallele Komponente $F_{G\parallel}$ der
Gewichtskraft, nicht die volle Gewichtskraft $F_g$, denn die
senkrechte Komponente $F_{G\perp}$ wird bereits von der Normalkraft
$F_N$ ausgeglichen. Nach $a$ auflösen:
\[
F_{G\parallel}=m\cdot a \quad\Longrightarrow\quad a=\frac{F_{G\parallel}}{m}
\]
Werte einsetzen:
\[
a=\frac{\SI{100}{\newton}}{\SI{30}{\kilogram}}\approx\SI{3.3}{\meter\per\second\squared}
\]
\end{solution}

\part[3] Bestimmen Sie mit $v^2=2\cdot a\cdot s$ die Endgeschwindigkeit
$v$ des Schlittens am Ende der Rampe.

\Antwortraum{2}

\begin{solution}
Nach $v$ auflösen:
\[
v^2=2\cdot a\cdot s \quad\Longrightarrow\quad v=\sqrt{2\cdot a\cdot s}
\]
Werte einsetzen:
\[
v=\sqrt{2\cdot\SI{3.3}{\meter\per\second\squared}\cdot\SI{12}{\meter}}
  \approx\SI{8.9}{\meter\per\second}
\]
\end{solution}

\part[2] Wie lange dauert die Hangabfahrt? Nutzen Sie dafür
$v=a\cdot t$ mit Ihren Werten für $a$ und $v$.

\Antwortraum{1.5}

\begin{solution}
Nach $t$ auflösen:
\[
v=a\cdot t \quad\Longrightarrow\quad t=\frac{v}{a}
\]
Werte einsetzen:
\[
t=\frac{\SI{8.9}{\meter\per\second}}{\SI{3.3}{\meter\per\second\squared}}
  \approx\SI{2.7}{\second}
\]
\end{solution}
\end{parts}
\end{taskbox}
\end{questions}

% =============================================================================
% Aufgabe 3: Normalkraft dreier Fahrzeuge (LZ2, andere Perspektive)
% =============================================================================
\newpage
\begin{questions}
\setcounter{question}{2}
\begin{taskbox}[Aufgabe 3: Fahrrad, Auto, Lastwagen]
\question[5] Ein Fahrrad, ein Auto und ein Lastwagen stehen alle still
auf derselben waagrechten Strasse.

\begin{parts}
\part[3] Ordnen Sie die drei Normalkräfte der Grösse nach (kleinste
zuerst), und begründen Sie kurz.

\Antwortraum{2}

\begin{solution}
Alle drei Fahrzeuge sind in Ruhe: Es gilt jeweils $F_N=F_g=m\cdot g$.
Da die Masse von Fahrrad, Auto und Lastwagen zunimmt, nimmt auch die
Normalkraft in dieser Reihenfolge zu:
$F_N(\text{Fahrrad}) < F_N(\text{Auto}) < F_N(\text{Lastwagen})$.
\end{solution}

\part[2] Alle drei Fahrzeuge parken nun stattdessen (mit angezogener
Bremse) auf einer abschüssigen Strasse. Ändert sich die Reihenfolge aus
(a)?

\Antwortraum{1.5}

\begin{solution}
Nein. Auf der geneigten Strasse wird bei allen drei Fahrzeugen derselbe
Anteil der Gewichtskraft als parallele Komponente abgezweigt, weil der
Neigungswinkel für alle drei gleich ist. Die Normalkraft wird deshalb
bei allen drei Fahrzeugen im gleichen Verhältnis kleiner, und an der
Reihenfolge ändert das nichts.
\end{solution}
\end{parts}
\end{taskbox}
\end{questions}

\end{document}
